% ----------------------------------------------------------
% Conclusão
% ----------------------------------------------------------
\chapter{Viabilidade Financeira}

A análise de viabilidade financeira tem como objetivo avaliar a sustentabilidade econômica do projeto Linha Vital considerando sua possível inserção no mercado de tecnologia assistiva e saúde digital.

Alinhado ao seu propósito social, o sistema foi concebido para disponibilizar gratuitamente suas funcionalidades ao usuário final, especialmente pessoas idosas, indivíduos com necessidades especiais e pessoas que vivem sozinhas. Dessa forma, a sustentabilidade financeira da solução foi projetada com base em um modelo de parcerias institucionais e corporativas, permitindo a manutenção da infraestrutura tecnológica sem custos diretos para os usuários.

A análise apresentada neste capítulo contempla os custos operacionais estimados para manutenção da aplicação, as possíveis fontes de receita e diferentes cenários de operação, permitindo avaliar a viabilidade econômica do projeto em condições distintas de adoção pelo mercado.

\section{Custos}

Os custos do projeto concentram-se principalmente na infraestrutura tecnológica necessária para execução da aplicação, armazenamento de dados e comunicação com os usuários. Também são considerados os custos relacionados ao licenciamento e publicação da aplicação em plataformas de distribuição digital.

Durante a fase de desenvolvimento, o trabalho realizado pelos integrantes da equipe foi conduzido de forma colaborativa e sem remuneração financeira direta. Dessa forma, não foram considerados custos de recursos humanos para a construção inicial do sistema, caracterizando uma contribuição baseada em capital intelectual dos próprios desenvolvedores.

A Tabela~\ref{tab_custos} apresenta a estimativa dos custos necessários para implantação e manutenção do sistema Linha Vital.

\begin{table}[H]
\caption{Estimativa de custos do projeto Linha Vital}
\label{tab_custos}
\centering
\footnotesize

\begin{tabular}{|p{4cm}|p{5cm}|p{3cm}|p{2.5cm}|}
\hline
\textbf{Classificação do Custo} & \textbf{Descrição do Recurso} & \textbf{Tipo de Gasto} & \textbf{Valor (R\$)} \\
\hline

Infraestrutura (Mensal) &
Servidor VPS (Oracle Cloud OCI) &
Fixo Mensal &
R\$ 11,00 \\
\hline

Infraestrutura (Mensal) &
Banco de Dados Relacional PostgreSQL &
Open-Source &
R\$ 0,00 \\
\hline

Infraestrutura (Mensal) &
Gateway de Disparo de SMS (Média Estimada) &
Variável Mensal &
R\$ 30,00 \\
\hline

Licenciamento (Único) &
Taxa de Desenvolvedor Google Play &
Investimento &
R\$ 130,00 \\
\hline

Recursos Humanos &
Desenvolvimento Técnico da Equipe &
Sweat Equity &
R\$ 0,00 \\
\hline

Total de Investimento Inicial Real &
Taxa Única de Publicação na Loja &
Inicial &
R\$ 130,00 \\
\hline

Total de Custo Fixo Operacional &
Sustentação Mensal da Tecnologia &
Recorrente &
R\$ 41,00/mês \\
\hline

\end{tabular}

\normalsize

\legend{Fonte: Elaborado pelos autores (2026).}
\end{table}

Cabe destacar que os valores apresentados na Tabela~\ref{tab_custos} representam os custos mínimos necessários para operação da versão inicial da solução (MVP). Para fins de análise de viabilidade financeira e projeção de operação comercial, considera-se um custo operacional mensal estimado de R\$ 320,00, contemplando a expansão da infraestrutura computacional, aumento do volume de comunicações, serviços complementares de suporte e recursos necessários para atender uma base crescente de usuários e instituições parceiras.
